\section{Technology Assessment}
\label{sec:technology}


Introduce in (sufficient) depth the key concepts and architecture of the chosen software technology. As part if this, you may consider using a running example to introduce the technology.

This part and other parts of the report probably needs to refer to
figures. Figure~\ref{fig:framework} from \cite{brown:96} just
illustrates how figure can be included in the report.

Entity Framework (EF) is an Object-Relational Mapper (ORM) developed by Microsoft for the .NET platform. An ORM supports the creation of applications focused on data management by abstracting the data storage and allows developers to focus on the business logic of their applications while interacting with underlying databases. This abstraction reduces the need to write extensive boilerplate code for database operations and provides a more intuitive object-oriented approach to data handling.

Entity Framework does this by dividing the application into three parts: the conceptual model, the storage model, and the mappings between them. The conceptual model represents the application's domain entities and relationships. It abstracts the data source, allowing developers to work with objects and relationships at a higher, business-oriented level. The storage model describes the logical structure of the database, including tables, columns, keys, and constraints. This model is specific to the underlying data source. Developers interact primarily with the conceptual model, making the application independent of the specific database technology. The Mapping specification language (MSL) defines the mappings between these two models, which specifies how entities, properties, and relationships in the conceptual model align with the database schema in the storage model. The conceptual model is defined using classes in code, and configurated to specify how these classes should be connected in the database.  

\begin{figure}[thb]
	\centering
	\includegraphics[scale=0.5]{figs/framework.png}
	\caption{Software technology evaluation framework.}
	\label{fig:framework}
\end{figure}

\subsection{Descriptive Modeling}

write where the technology comes from, its history, its context and what problem it solves.
Consider drawing a graph like in \cite{brown:96}.

\textbf{Technology genealogy}
The first version of Entity Framework was released in 2008, as part of .NET Framework 3.5 SP1 and Visual Studio 2008 SP1, by Microsoft. Prior to EF, developers relied on ADO.NET for data access within .NET applications. Entity Framework (EF) was developed as an answer to two different objectives, to raise the level of abstraction when working with modelling of entities, and data engines to store and retrieve data. EF enables developers to work with data in the form of domain-specific objects and properties, such as customers and customer addresses, without having to concern themselves with the underlying database tables and columns where this data is stored. With EF, developers can work at a higher level of abstraction when they deal with the data and can create and maintain data-oriented applications with less code than in traditional applications [1]. Its first version was widely criticized, even attracting a “vote of no confidence” signed by at least one thousand developers [3]. The second version, EF 4.0, was as part of .NET 4.0 in 2010 and addresses many of the criticisms made of version 1, such as lazy loading, testability improvements, foreign key mapping, support for Plain old CLR object (POCO) which is a simple object created in .NET, often used in opposition to complex or specialized objects that object-relational mapping framework often require [5].  With the release of version 6.0 in 2013, EF became an open source project licensed under Apache License v2. With this version, EF was also moved out of the band from the .NET Framework, which helps in accelerating the pace of development and delivery of new features [4]. 


\subsection{Experiment Design}
\label{subsec:experiment_design}
Keeping in line with our evaluation framework \cite{brown:96} we are to first do a comparative feature analysis before we further our attention towards creating hypotheses and experiment design.
To start of our experiment design we create a comparative feature analysis.

\subsubsection{Comparative feature analysis}
To conduct our comparative feature analysis we created Table~\ref{tab:jpa_vs_ef_vs_no_orm}, in this table we have looked at the official documentation of the Entity Framework \cite{ef_getting_started} and the official documentation of JPA \cite{spring_jpa}.
\begin{table}[H]
    \centering
    \begin{tabularx}{\textwidth}{|l|X|X|X|} % Define column widths with borders
        \hline
        \textbf{Feature} & \textbf{JPA} & \textbf{Entity Framework} & \textbf{Not Using ORM} \\ \hline
        Schema Abstraction & Relies on annotations and configuration (e.g., Hibernate supports @Table, @Column). & Strong tooling with migrations (Add-Migration, Update-Database commands). & Schema directly defined and managed in SQL scripts or database-specific tools. \\ \hline
        Schema Evolution Tools & Schema auto-generation and validation with support for tools like Flyway and Liquibase. & Built-in migration system for versioning and rolling back schema changes. & Requires manual management of schema changes and evolution. \\ \hline
        Query Abstraction & JPQL for writing abstract, database-independent queries. & LINQ provides strongly-typed, language-integrated query capabilities. & Direct SQL queries must be written and maintained, often specific to the database. \\ \hline
        CRUD Boilerplate Reduction & Automatic through frameworks like Spring Data JPA. & Automatic using EF Core Repositories and DbContext for basic CRUD operations. & No automation; full CRUD operations must be manually implemented in SQL or procedural code. \\ \hline
        Ease of Refactoring & Code annotations simplify model changes but may require manual migration scripting. & Strong support through auto-generated migration files and IDE integration. & Refactoring requires careful updates to SQL scripts and application code to avoid inconsistencies. \\ \hline
        Cross-Platform Support & Platform-agnostic via the JVM, supports multiple database engines. & Cross-platform with .NET Core but optimized for Windows and SQL Server. & Fully dependent on the database and programming environment used. \\ \hline
        Code Complexity Reduction & Annotations reduce code overhead for mappings but require verbose configuration files. & Strongly-typed models and DbContext minimize configuration and boilerplate code. & High complexity due to explicit SQL queries and manual schema synchronization. \\ \hline
    \end{tabularx}
    \caption{Comparison of JPA, Entity Framework, and Not Using ORM}
    \label{tab:jpa_vs_ef_vs_no_orm}
\end{table}


As shown in Table \ref{tab:jpa_vs_ef_vs_no_orm} we see that there are some clear similarities and some differences. 
As a proper first step we wanted to focus our attention on ORM´s contribution to a project like our own.
This helped us generate the hypotheses found in the next section.

\subsubsection{Hypotheses}
Looking at the ORM's we could see that what we could call programmer comfort is an important part of what the ORM's contribution to the project. 
Another important part looks to be the reduction of repetitive or "boilerplate" code.
With this in mind we created the following hypotheses:

\textbf{Hypotheses one} \newline
The abstraction provided by an ORM simplifies maintenance and refactoring of database schemas. \\

\textbf{Hypotheses two} \newline
ORM shortens development time by reducing unnecessary boilerplate code. \\

We believe these hypotheses would give an interesting insight regarding the possible benefits of utilizing an ORM.

\subsubsection{Experiment design}
As the name of the subchapter \ref{subsec:experiment_design} hints to we are to design an experiment. 
To test our hypotheses we want to create some experiments that are possible to recreate and that are measurable. \\

\textbf{Experiment for hypotheses one} \\
To test if an ORM simplifies maintenance and refactoring of database schemas we need to simulate schema evolution scenarios in ORMs and without ORM.
To be certain that the experiment is measuring the correct parameters we want to make sure that the database technology is the same and in our case that means using PostgreSQL, and we also want to make the same changes to the table in all experiments.
To measure the differences we want to look at the following: 
\begin{itemize}
    \item Measure the time it takes to apply schema changes.
    \item Developer effort (lines of code modified, complexity of changes)
    \item Risk of introducing errors (While timing the schema change time also record any breaking changes)
\end{itemize}


\textbf{Experiment for hypotheses two}
To test if an ORM shortens development time by reducing unnecessary boilerplate code we need to build a CRUD application.
To be certain that the experiment is measuring the correct parameters we need to have some controlled variables like all applications needs to have the same functional requirement.
To measure the differences we want to look at the following:
\begin{itemize}
	\item Development time.
	\item Code size
	\item Complexity and readability (measured by developer feedback)
\end{itemize}

\subsection{Experiment Evaluation}
The experiment evaluation face we are sectioning into setup, methodology, results, discussion and limitations.\\
The experiments for both hypotheses as outlined in the previous subsection \ref{subsec:experiment_design} the fist hypotheses is dependent on a completed project and the second is largely dependent on creating a project.
These dependencies give us the opportunity to combine the valuation experiments into one experiment.
The experiment evaluates how ORM frameworks (JPA and EF) \textbf{impact development time}  and \textbf{schema maintenance}.\\ 
The hypothesis states that ORM abstraction:
\begin{enumerate}
	\item ORM abstraction simplifies schema maintenance compared to raw SQL.
	\item ORM reduces development time by minimizing repetitive code.
\end{enumerate}

\subsubsection{Experiment setup}
We want to assess CRUD implementation and schema maintenance efficiency for JPA, EF Core and raw SQL.

\textbf{Environment} \\
We want the environments to be as similar as possible, but some differences are present.
\begin{itemize}
	\item PostgreSQL is used for all applications.
	\item For EF Core Visual Studio 2022 is used.
	\item For JPA IntelliJ IDEA is used.
\end{itemize}

\textbf{Database model}
The database schema needs to be the same for all applications.
\begin{verbatim}
	Poll(ID, Question, PublishedAt, ValidUntil, VoteOptions)
	VoteOption(ID. Caption, PresentationOrder)
\end{verbatim}

\textbf{Tasks} \\
The experiment is to be conducted by implementing CRUD operations for Poll and VoteOption.
When the CRUD operations are implemented \textbf{Title} is to be added to \textbf{Poll}.

\subsubsection{Methodology}
Due to constraints of time the experiment involves only one developer.

\textbf{Metrics} \\
We measure the following:
\begin{itemize}
	\item \textbf{Development Time}: Time to complete each task.
	\item \textbf{Lines of Code (LOC)}: To quantify boilerplate reduction.
	\item \textbf{Error Rate}: Number of runtime errors encountered during task completion.
\end{itemize}

\textbf{Process} \\
\begin{itemize}
	\item The developer is timed while implementing CRUD operations and performing the schema modification.
	\item Errors are logged, and corrections incur time penalties.
	\item Tasks are executed in a consistent environment to minimize external influences.
\end{itemize}

\subsubsection{Results}
\textbf{Development Time} \\
\begin{tabular}{|l|c|c|c|}
\hline
\textbf{Task} & \textbf{JPA (Hibernate)} & \textbf{EF Core} & \textbf{Raw SQL} \\ \hline
CRUD Operations & 45 min & 30 min & 60 min \\ \hline
Schema Modification & 25 min & 5 min & 40 min \\ \hline
\end{tabular}

\textbf{Lines of Code (LOC)} \\
\begin{tabular}{|l|c|c|c|}
\hline
\textbf{Task} & \textbf{JPA (LOC)} & \textbf{EF Core (LOC)} & \textbf{Raw SQL (LOC)} \\ \hline
CRUD Operations & 120 & 90 & 200 \\ \hline
Schema Modification & 30 & 20 & 60 \\ \hline
\end{tabular}

\textbf{Error Rate} \\
\begin{tabular}{|l|c|c|c|}
\hline
\textbf{Task} & \textbf{JPA} & \textbf{EF Core} & \textbf{Raw SQL} \\ \hline
CRUD Operations & 2 & 1 & 4 \\ \hline
Schema Modification & 1 & 0 & 3 \\ \hline
\end{tabular}

\subsubsection{Discussion}
\textbf{Development Time} \\
EF Core demonstrated the shortest development time due to its integrated migration tools and LINQ support, which makes schema updates and query construction easier. JPA was similar in development time, but its annotation-based configuration added up to more time spent programming. Raw SQL required the most time, as every query and schema change had to be manually implemented.

\textbf{Lines of Code (LOC)} \\
EF Core had the least LOC due to its DbContext abstraction, reducing boilerplate code to what seems like a minimum. JPA required more lines due to annotations. Raw SQL added the most LOC due to requiring explicit query and schema manipulation logic.

\textbf{Error Rate} \\
EF Core had the lowest error rate, benefiting from compile-time checks for LINQ queries and Visual Studio debugging tools. JPA had slightly more runtime errors, as JPQL queries are validated at runtime. Raw SQL had the highest error rate due to manual query construction, which increased typos and logic errors.

\subsubsection{Conclusion}
The experiment shows that ORMs (EF Core and JPA) significantly reduce development time and boilerplate code compared to raw SQL. EF Core emerged as the most efficient ORM due to its tooling and abstraction, but JPA performed well in cross-platform scenarios. Raw SQL offers maximum control and may be preferable for performance-critical applications.

